\chapter{Laryngitis}

\section*{Definition}
Laryngitis is inflammation of the larynx (voice box), most commonly caused by viral upper respiratory infections. It causes hoarseness or loss of voice. Acute laryngitis lasts less than 3 weeks, while chronic laryngitis persists longer.

\section*{What Happens in Your Body}
Viral infection of the upper respiratory tract causes edema and inflammation of the vocal cords. The vocal cords become red, swollen, and less pliable, disrupting normal vibration and producing hoarseness. Overuse of the voice during inflammation can exacerbate the condition. Chronic laryngitis involves persistent irritation from reflux, smoking, or environmental factors.

\section*{Causes}
\begin{itemize}
\item Viral upper respiratory infections (most common)
\item Vocal strain from yelling, singing, or prolonged speaking
\item Gastroesophageal reflux disease (laryngopharyngeal reflux)
\item Bacterial infections (less common)
\item Allergies and postnasal drip
\item Inhaled irritants (cigarette smoke, pollution)
\item Chronic sinusitis
\item Fungal laryngitis (immunosuppressed people)
\item Autoimmune conditions
\item Intubation or surgery involving the airway
\end{itemize}

\section*{When to See a Doctor}
\begin{itemize}
\item Hoarseness or raspy voice
\item Weak voice or loss of voice (aphonia)
\item Sore throat or throat discomfort
\item Dry, barking cough
\item Feeling of a lump in the throat
\item Frequent throat clearing
\item Pain with speaking or swallowing
\item Low-grade fever
  \item Hoarseness lasts more than 2--3 weeks
\end{itemize}

\section*{Related Symptoms}
\begin{itemize}
\item Sore throat
\item Cough
\item Hoarseness
\item Postnasal drip
\item GERD (heartburn)
\item Upper respiratory infection (common cold)
\item Tonsillitis
\item Pharyngitis
\end{itemize}

\section*{Self-Care/Home Management}
\begin{itemize}
\item Voice rest: avoid speaking, whispering, or shouting
\item Drink warm, soothing fluids (honey and lemon tea, broth)
\item Use a humidifier or inhale steam
\item Avoid smoking and secondhand smoke
\item Avoid alcohol and caffeine (can dehydrate)
\item Gargle with warm salt water
\item Use OTC pain relievers for discomfort
\item Do NOT whisper (strains the vocal cords more than normal speech)
\end{itemize}

\section*{How Your Doctor Will Evaluate You}
\begin{itemize}
\item Clinical history and symptom pattern
\item Laryngoscopy (mirror or flexible endoscopy) to visualize vocal cords
\item Referral to otolaryngologist for chronic or atypical cases
\item Laryngeal stroboscopy for detailed vocal cord assessment
\item pH monitoring if reflux laryngitis suspected
\end{itemize}

\section*{Treatment Options}
\begin{itemize}
\item Voice rest and hydration
\item Treat underlying cause (antibiotics for bacterial, PPI for reflux)
\item Corticosteroids for severe inflammation or airway compromise
\item Speech therapy for vocal cord dysfunction or chronic hoarseness
\item Surgery for vocal cord lesions (polyps, nodules) if present
\end{itemize}
